\documentclass[11pt]{article}

\usepackage[margin=1in]{geometry}
\usepackage{amsmath, amssymb, amsthm}
\usepackage{booktabs}
\usepackage{enumitem}
\usepackage{xcolor}
\usepackage{bm}
\usepackage[colorlinks=true, linkcolor=blue!60!black, urlcolor=blue!60!black,
            citecolor=blue!60!black]{hyperref}
\usepackage{listings}

\lstset{
  basicstyle=\ttfamily\small, breaklines=true, columns=fullflexible,
  keepspaces=true, frame=single, framesep=4pt,
  backgroundcolor=\color{gray!7}, xleftmargin=6pt,
}

\newcommand{\code}[1]{\texttt{#1}}
\newcommand{\R}{\mathbb{R}}
\newcommand{\E}{\mathbb{E}}
\newcommand{\Cov}{\mathrm{Cov}}
\newcommand{\bx}{\bm{x}}
\newcommand{\by}{\bm{y}}
\newcommand{\bm}[1]{\boldsymbol{#1}}
\newcommand{\bmu}{\bm{\mu}}
\newcommand{\bSig}{\bm{\Sigma}}
\newcommand{\half}{\tfrac{1}{2}}
\DeclareMathOperator{\diag}{diag}

\newtheorem{lemma}{Lemma}
\newtheorem{proposition}{Proposition}

\title{The Sigma-Point (Unscented) Forecast for \texttt{ORJointUpdate}:\\
Seven Points for the Undisturbed and Every Disturbance Class}
\author{Notes on \texttt{OR\_joint\_update.Rmd}}
\date{15 July 2026}

\begin{document}
\maketitle

\begin{abstract}
Profiling showed that the joint filter's cost is essentially \emph{linear in the
ensemble size} $J$: the VSEM pushforward, the disturbance draws, and the analysis
resample all run once per ensemble member. This note develops the alternative that
removes the $J$ axis entirely. The forecast step needs only the \emph{mean and
$3\times3$ covariance} of the pushforward for each mixture label, so we replace the
$J\!\approx\!100$ random members by $2n+1=7$ deterministic \emph{sigma points} (for the
$n=3$ state) and reconstruct those two moments with the unscented transform. The central
observation is that all $D$ labels --- undisturbed and each of the \code{dist\$n}
disturbance classes --- are pushforwards of the \emph{same} input Gaussian through maps
that \emph{share the same expensive VSEM stage}. So VSEM is evaluated at the seven points
\emph{once}; each label then reuses those seven propagated points and applies only its own
cheap, closed-form post-map. The disturbance operator's intrinsic randomness is folded in
analytically by the law of total variance, using conditional moments we derive in closed
form. Net effect: $7$ VSEM calls per component per step instead of $J$, no Monte-Carlo
subsetting noise, and the ensemble representation (and its resample) disappears.
\end{abstract}

\tableofcontents

\section{Why this helps, in one paragraph}
\label{sec:why}

The Rao--Blackwellization of the precision (companion note
\code{claude-speedups-7.15.26.tex}, Idea~1) removed the quadrature factor $G$ from the
VSEM cost. Profiling the result revealed that VSEM was never the dominant term: the run
time scales \emph{linearly in $J$}, dominated by (i) the per-member disturbance draw
\code{disturbance.p}, (ii) the Kalman/collapse arithmetic, and (iii) VSEM itself, each
paid $\bigO(J)$ times per component. The forecast only ever \emph{uses} a mean and a
$3\times3$ covariance per label (see \code{vsem\_pushforward}, which returns
\code{colMeans} / \code{var} of the pushforward). Estimating two moments of a smooth map
does not require $100$ random samples --- $7$ well-placed deterministic points do it to
second order. That is the unscented transform, and it collapses the $J$ axis to a
constant $7$.

\section{Setup and notation}
\label{sec:setup}

\paragraph{State.} The model state is the $n=3$ carbon-pool vector
\[
\bx = (x_1, x_2, x_3) = (C_v,\; C_s,\; C_R)
   = (\text{foliage biomass},\; \text{soil carbon},\; \text{wood/AGB}),
\]
and the observation operator picks out the wood pool, $H=(0,0,1)$, so $x_3$ is the
observed quantity (\code{OR\_joint\_update.Rmd}).

\paragraph{Input to the forecast.} At each step, every mixture component $c$ carries a
Gaussian analysis $\mathcal N(\bm m, \bm P)$ over $\bx$ (in the current code this is the
$\Pi_D$ result \code{mu.k}/\code{Sig.k}, from which the ensemble \code{ens} is
resampled). \emph{All $D$ labels forecast from this same $\mathcal N(\bm m,\bm P)$}; they
differ only in the map applied.

\paragraph{The two maps.} Let $f_v:\R^3\!\to\!\R^3$ be the deterministic VSEM
integration over the year's PAR (the expensive stage). The labels are:
\begin{itemize}[nosep]
  \item \textbf{Undisturbed} ($k=1$): $\bx \mapsto f_v(\bx)$, then additive process error
        $\mathcal N(\bm 0, Q_g)$ with $Q_g[3,3]=1/\tau_g$ folded in analytically (this is
        the node-specific piece from the RB note; it does not touch VSEM).
  \item \textbf{Disturbance class $d$} ($k=d+1$, $d=1..\code{dist\$n}$):
        $\bx \mapsto g_d\!\big(f_v(\bx)\big)$, where $g_d$ is the \emph{stochastic}
        disturbance operator \code{disturbance.p} with class parameters
        $(\bmu_0^{(d)}, V_0^{(d)})$ and soil-allocation fraction $s=\code{alloc.soil}$.
\end{itemize}
The critical structural fact, which the whole speedup rests on:
\begin{center}
\fbox{\parbox{0.92\linewidth}{\centering
$f_v$ is common to \emph{all} labels. Only the cheap post-map differs
($\text{identity}+Q_g$ for undisturbed, $g_d$ for class $d$).}}
\end{center}

\paragraph{Notation.}
\begin{center}
\begin{tabular}{@{}ll@{}}
\toprule
$n=3$ & state dimension \\
$2n+1=7$ & number of sigma points \\
$\bm m,\ \bm P$ & mean / covariance of the component's input Gaussian \\
$\mathcal X_i,\ i=0..2n$ & sigma points \\
$w_i^m,\ w_i^c$ & mean / covariance reconstruction weights \\
$\mathcal Y_i=f_v(\mathcal X_i)$ & VSEM-propagated sigma points (\emph{shared across labels}) \\
$\mu_d(\by),\ \Sigma_d(\by)$ & conditional mean/cov of $g_d$ given its input $\by$ \\
$s$ & \code{alloc.soil}, soil allocation fraction \\
\bottomrule
\end{tabular}
\end{center}

\section{The unscented transform}
\label{sec:ut}

We want the mean and covariance of $\by=\phi(\bx)$ for a smooth map $\phi$, when
$\bx\sim\mathcal N(\bm m,\bm P)$. The unscented transform
\cite{julier2004,vandermerwe2004,sarkka2013} chooses $2n+1$ deterministic points that
match $\bm m$ and $\bm P$ exactly, pushes them through $\phi$, and reconstructs the output
moments as weighted sample statistics.

\paragraph{Sigma points.} With scaling parameter $\lambda=\alpha^2(n+\kappa)-n$,
\begin{equation}
\label{eq:sigma}
\mathcal X_0=\bm m,\qquad
\mathcal X_i=\bm m+\big[\sqrt{(n+\lambda)\bm P}\,\big]_i,\qquad
\mathcal X_{i+n}=\bm m-\big[\sqrt{(n+\lambda)\bm P}\,\big]_i,\quad i=1..n,
\end{equation}
where $\big[\sqrt{M}\,\big]_i$ is the $i$-th column of a matrix square root of $M$
(a Cholesky factor $L$ with $LL^\top=M$ is the usual choice). For $n=3$ this is
$\mathcal X_0$ plus $3$ ``$+$'' points and $3$ ``$-$'' points, i.e.\ \textbf{seven points},
one at the mean and a symmetric pair along each principal direction of $\bm P$.

\paragraph{Weights.}
\begin{equation}
\label{eq:weights}
w_0^m=\frac{\lambda}{n+\lambda},\quad
w_0^c=\frac{\lambda}{n+\lambda}+(1-\alpha^2+\beta),\quad
w_i^m=w_i^c=\frac{1}{2(n+\lambda)}\ \ (i=1..2n).
\end{equation}
Standard Gaussian choices are $\alpha=10^{-3}$ (small spread), $\kappa=0$, $\beta=2$
(optimal for Gaussians). The weights sum to one, $\sum_i w_i^m=1$.

\paragraph{Reconstruction.} With $\mathcal Y_i=\phi(\mathcal X_i)$,
\begin{equation}
\label{eq:recon}
\hat{\bm m}=\sum_{i=0}^{2n} w_i^m\,\mathcal Y_i,\qquad
\hat{\bm P}=\sum_{i=0}^{2n} w_i^c\,(\mathcal Y_i-\hat{\bm m})(\mathcal Y_i-\hat{\bm m})^\top .
\end{equation}

\begin{proposition}[Exactness for affine maps]
\label{prop:affine}
If $\phi(\bx)=A\bx+\bm b$ then \eqref{eq:recon} returns $\hat{\bm m}=A\bm m+\bm b$ and
$\hat{\bm P}=A\bm P A^\top$ exactly, for any valid $(\alpha,\beta,\kappa)$.
\end{proposition}
\noindent For nonlinear $\phi$ the transform is exact to second order in the Taylor
expansion (third order in mean for symmetric points), which is why $7$ deterministic
points beat $100$ random ones for estimating two moments: the random estimator carries
$\bigO(J^{-1/2})$ Monte-Carlo error in $\hat{\bm P}$, while the UT has none for the
affine part and small bias for the curvature. A short proof of
Proposition~\ref{prop:affine} is in Appendix~\ref{app:affine}.

\section{Seven points, shared across all labels}
\label{sec:shared}

This section answers the question directly: \emph{how do seven points serve the
undisturbed label and every disturbance class at once?}

\paragraph{Step 1 --- build the seven points from the input Gaussian.} Use
\eqref{eq:sigma} on the component's $\mathcal N(\bm m,\bm P)$. This is pure linear
algebra (one Cholesky of a $3\times3$), independent of any label.

\paragraph{Step 2 --- push them through VSEM \emph{once}.} Evaluate
\begin{equation}
\label{eq:propagate}
\mathcal Y_i=f_v(\mathcal X_i),\qquad i=0,1,\dots,6,
\end{equation}
i.e.\ \textbf{seven VSEM integrations total} (versus $J$ today). Because $f_v$ is the
common stage of every label's map, the propagated set $\{\mathcal Y_i\}$ is all the
expensive information any label needs. \emph{Nothing below re-runs VSEM.}

\paragraph{Step 3a --- undisturbed moments.} Apply \eqref{eq:recon} to
$\{\mathcal Y_i\}$ and add the process error analytically:
\begin{equation}
\label{eq:undist}
\bm m_1=\sum_i w_i^m\,\mathcal Y_i,\qquad
\bm P_1=\underbrace{\sum_i w_i^c\,(\mathcal Y_i-\bm m_1)(\mathcal Y_i-\bm m_1)^\top}_{\text{state pushforward spread}}
        \;+\;\underbrace{Q_g}_{\text{process error}} .
\end{equation}
This mirrors exactly what the current code does with \code{colMeans}/\code{var} plus
\code{+ Qg}; the process-error part $Q_g[3,3]=1/\tau_g$ is still isolated, so the NIS
decomposition (\code{nis.q33}) is preserved unchanged.

\paragraph{Step 3b --- each disturbance class.} Reuse the \emph{same}
$\{\mathcal Y_i\}$; apply the class-$d$ operator. Because $g_d$ is stochastic we do
\emph{not} sample it --- we use its closed-form conditional moments $\mu_d,\Sigma_d$
(derived in \S\ref{sec:disturb}) and the law of total variance:
\begin{align}
\bm m_{d+1}&=\sum_i w_i^m\,\mu_d(\mathcal Y_i), \label{eq:distmean}\\
\bm P_{d+1}&=\underbrace{\sum_i w_i^c\,\big(\mu_d(\mathcal Y_i)-\bm m_{d+1}\big)\big(\mu_d(\mathcal Y_i)-\bm m_{d+1}\big)^\top}_{\text{between-point: state uncertainty through }g_d}
            \;+\;\underbrace{\sum_i w_i^m\,\Sigma_d(\mathcal Y_i)}_{\text{within: intrinsic disturbance noise}}. \label{eq:distcov}
\end{align}
Equations \eqref{eq:distmean}--\eqref{eq:distcov} are the multi-point law of total
variance: $\Cov(g_d)=\Cov\big(\E[g_d\mid \by]\big)+\E\big[\Cov(g_d\mid \by)\big]$, with
the outer $\by$-distribution represented by the seven propagated points and their
weights. The first (``between'') term is the disturbance-model mean pushed through the
state spread; the second (``within'') term is the disturbance draw's own variance,
averaged over the points --- precisely the piece the current code gets by \emph{sampling}
\code{disturbance.p} per member.

\paragraph{Why this is legitimate --- and better.} In the ensemble code the multinomial
split (\code{rmultinom}) assigns each member to one label, so each label's moments are
estimated from a \emph{random subset} of members. That split is only a Monte-Carlo device
to ``apply a different map to different draws''; it injects subsetting noise and makes a
label's covariance depend on how many members happened to land in it. The label
\emph{weights} $\beta_k$ come from the disturbance prior in the main loop, \emph{not} from
the split counts. The UT version discards the split: it applies \emph{every} label's map
to \emph{all seven} propagated points, deterministically. This is strictly more accurate
(no subsetting noise, no small-count covariance blow-ups) and is what lets one shared
seven-point set serve all $D$ labels.

\paragraph{Cost.} Per component per step: one $3\times3$ Cholesky, $7$ VSEM integrations,
and $1+\code{dist\$n}$ cheap analytic post-maps (each $\bigO(7)$ tiny operations). VSEM
calls drop from $J$ to $7$; the disturbance draws ($\bigO(J)$ \code{rmvnorm} calls, the
profiling hotspot) vanish entirely, replaced by closed-form $\mu_d,\Sigma_d$; and since
we now carry $(\bm m,\bm P)$ directly, the per-step analysis \code{rmvnorm(J,\dots)}
resample also disappears.

\section{The disturbance operator's conditional moments}
\label{sec:disturb}

We now derive $\mu_d(\by)$ and $\Sigma_d(\by)$ used in
\eqref{eq:distmean}--\eqref{eq:distcov}. Recall \code{disturbance.p} acting on a
post-VSEM state $\by=(y_1,y_2,y_3)=(\text{leaf},\text{soil},\text{stem})$
(\code{disturbance.R}). It treats the biomass pair $\bm u=(y_1,y_3)$ (indices
\code{c(1,3)}), draws a post-disturbance biomass
\begin{equation}
\label{eq:draw}
\bm b=(b_1,b_3)\sim\mathcal N\!\big(\bmu_0\odot\bm u,\; D_0(\bm u)\big),
\qquad D_0(\bm u)=\big(\bm u\bm u^\top\big)\odot V_0,
\end{equation}
(the code's \code{tcrossprod(old.biomass)*V0}), then sets the new state
\begin{equation}
\label{eq:newstate}
o_1=b_1,\qquad o_3=b_3,\qquad o_2=y_2+s\big[(y_1-b_1)+(y_3-b_3)\big],
\end{equation}
i.e.\ the removed biomass $\,r=(y_1-b_1)+(y_3-b_3)\,$ has a fraction $s$ deposited in
soil. ($\bmu_0=\bmu_0^{(d)}$, $V_0=V_0^{(d)}$; $\odot$ is elementwise.)

\paragraph{$g_d$ is affine in the Gaussian draw.} Holding $\by$ fixed, \eqref{eq:newstate}
is linear in $\bm b$:
\begin{equation}
\label{eq:affine}
\bm o = A\,\bm b + \bm c,\qquad
A=\begin{pmatrix} 1 & 0\\ -s & -s\\ 0 & 1 \end{pmatrix},\qquad
\bm c=\begin{pmatrix} 0\\ y_2+s(y_1+y_3)\\ 0\end{pmatrix},
\end{equation}
with $\bm b$ Gaussian from \eqref{eq:draw}. Hence $\bm o\mid\by$ is Gaussian and its
moments are exact (no approximation in the disturbance operator itself; the only UT
approximation is in the $f_v$ stage of \S\ref{sec:shared}).

\paragraph{Conditional mean.} With $\E[\bm b\mid\by]=\bmu_0\odot\bm u
=(\mu_{0,1}y_1,\ \mu_{0,2}y_3)$,
\begin{equation}
\label{eq:mu_d}
\mu_d(\by)=A(\bmu_0\odot\bm u)+\bm c=
\begin{pmatrix}
\mu_{0,1}\,y_1\\[2pt]
y_2+s\big[(1-\mu_{0,1})y_1+(1-\mu_{0,2})y_3\big]\\[2pt]
\mu_{0,2}\,y_3
\end{pmatrix}.
\end{equation}
This is exactly the deterministic disturbance: retain a fraction $\mu_{0,\cdot}$ of leaf
and stem, and route fraction $s$ of what was removed into soil.

\paragraph{Conditional covariance.} With $V_0=\diag(v_1,v_2)$ (as in the code), the draw
covariance \eqref{eq:draw} is diagonal,
$D_0(\bm u)=\diag\!\big(y_1^2 v_1,\; y_3^2 v_2\big)=:\diag(d_1,d_2)$. Then
$\Sigma_d(\by)=A\,D_0\,A^\top$ evaluates to
\begin{equation}
\label{eq:Sig_d}
\Sigma_d(\by)=
\begin{pmatrix}
d_1 & -s\,d_1 & 0\\
-s\,d_1 & s^2(d_1+d_2) & -s\,d_2\\
0 & -s\,d_2 & d_2
\end{pmatrix},
\qquad d_1=y_1^2\,v_1,\ \ d_2=y_3^2\,v_2 .
\end{equation}
The structure is intuitive: leaf ($o_1$) and stem ($o_3$) inherit the draw variances
$d_1,d_2$; soil ($o_2$) absorbs fraction $s$ of both removals, giving variance
$s^2(d_1+d_2)$ and negative covariances with leaf and stem (more retained $\Rightarrow$
less to soil). Leaf and stem are uncorrelated because $V_0$ is diagonal; a full $V_0$
would add the cross term $-\,(\cdot)$ and an $o_1$--$o_3$ entry, obtained the same way
from $A D_0 A^\top$.

\paragraph{Two practical caveats.}
\begin{enumerate}[nosep]
  \item \textbf{Retention clamp.} \code{disturbance.p} constrains
        $b\le \bm u$ (no biomass gain) and applies \code{tobit} ($\ge 0$). The affine
        moments \eqref{eq:mu_d}--\eqref{eq:Sig_d} ignore these truncations. Since
        $\mu_{0,\cdot}<1$ the mean draw is safely below $\bm u$ and the clamp bites only
        in the tail; the resulting bias in the moments is small. If needed it can be
        corrected with truncated-Gaussian moment formulas, but for a forecast-covariance
        estimate the uncorrected form is normally adequate.
  \item \textbf{Backup branch.} The current code already falls back to fixed prior
        disturbance moments when a class has too few members to form a covariance
        (\code{OR\_joint\_update.Rmd}). Under the UT there is no member count, so this
        branch simply disappears --- $\mu_d,\Sigma_d$ are always well defined from
        \eqref{eq:mu_d}--\eqref{eq:Sig_d}.
\end{enumerate}

\section{Algorithm}
\label{sec:algo}

\begin{lstlisting}[caption={UT replacement for \code{vsem\_pushforward}, per component.}]
ut_pushforward <- function(m, P, tsel) {          # m: 3-vec, P: 3x3
  n <- 3; alpha <- 1e-3; beta <- 2; kappa <- 0
  lam <- alpha^2 * (n + kappa) - n
  L   <- t(chol((n + lam) * (P + diag(1e-9, 3)))) # columns = sqrt((n+lam)P)
  X   <- cbind(m, m + L, m - L)                   # 3 x 7 sigma points
  wm  <- c(lam/(n+lam), rep(1/(2*(n+lam)), 2*n))
  wc  <- wm; wc[1] <- wc[1] + (1 - alpha^2 + beta)

  Y <- apply(X, 2, function(x) VSEM(setTheta(x), PAR = PAR[tsel])[365, 2:4])  # 3 x 7
  Y <- pmax(Y, 1e-6)                              # positivity (tobit) guard

  ## --- undisturbed: reconstruct moments, add Qg analytically ---
  m1 <- as.numeric(Y %*% wm)
  P1 <- Reduce(`+`, lapply(1:7, function(i){d<-Y[,i]-m1; wc[i]*tcrossprod(d)})) + Qg
  cmu <- list(m1); cSig <- list(P1)

  ## --- each disturbance class d: reuse the SAME propagated Y ---
  for (d in 1:dist$n) {
    mu0 <- dist$vals[[d]]$mu0; V0 <- dist$vals[[d]]$V0; s <- dist$alloc.soil
    MU <- sapply(1:7, function(i) mu_d(Y[,i], mu0, s))          # eq (mu_d), 3 x 7
    md <- as.numeric(MU %*% wm)
    between <- Reduce(`+`, lapply(1:7, function(i){e<-MU[,i]-md; wc[i]*tcrossprod(e)}))
    within  <- Reduce(`+`, lapply(1:7, function(i) wm[i]*Sig_d(Y[,i], V0, s)))  # eq (Sig_d)
    cmu[[d+1]] <- md; cSig[[d+1]] <- between + within
  }
  list(mu = cmu, Sig = cSig)                       # same interface as before (no q33 needed:
}                                                  # undisturbed Sig = spread + Qg as in eq (undist))
\end{lstlisting}

\noindent The return value has the identical shape (\code{mu}, \code{Sig} lists of length
$D$, undisturbed covariance $=$ spread $+Q_g$) that the main loop already consumes, so the
Kalman update, $\Pi_K/\Pi_D$ collapses, evidence accumulation, and NIS code are unchanged.
The single behavioural difference is upstream: the component now carries $(\bm m,\bm P)$
and no ensemble, so the per-step resample \code{rmvnorm(J,\dots)} is deleted and $(\bm
m,\bm P)$ is set directly from the $\Pi_D$ collapse.

\section{Accuracy, failure modes, and a fallback}
\label{sec:accuracy}

\begin{itemize}[nosep]
  \item \textbf{When the UT is excellent.} $f_v$ (VSEM annual pool dynamics) is smooth and
        close to affine over the analysis-covariance scale, so the two-moment
        reconstruction is very accurate; the disturbance operator is handled
        \emph{exactly} (\S\ref{sec:disturb}), not approximated.
  \item \textbf{When to be careful.} The UT tracks only the first two moments. If a
        component's forecast is genuinely multimodal or heavy-tailed \emph{within a single
        label}, the Gaussian summary loses that --- but the filter already summarizes each
        label by a single Gaussian (\code{mm}), so no information is lost relative to the
        current design. Strong nonlinearity of $f_v$ across a very wide $\bm P$ can bias
        $\hat{\bm P}$; the scaling $\alpha$ controls the point spread and can be tuned.
  \item \textbf{Positivity.} Sigma points can stray to nonphysical (negative) pools in the
        tails; clamp with \code{tobit}/\code{pmax} before and after VSEM, as the ensemble
        code already does.
  \item \textbf{Near-singular $\bm P$.} Add a small jitter $\varepsilon I$ before the
        Cholesky (shown in the listing), matching the existing
        \code{(Ck+t(Ck))/2 + diag(1e-6,3)} guard.
  \item \textbf{Validation / fallback.} Because the interface is unchanged, one can run
        UT and the $J$-member ensemble side by side on a few sites and compare
        \code{probD}, the analysis mean/CI, and the NIS calibration. Keep the ensemble
        path behind a flag as a fallback for any site where the two diverge beyond
        Monte-Carlo error.
\end{itemize}

\appendix
\section{Proof of Proposition~\ref{prop:affine}}
\label{app:affine}

Let $\phi(\bx)=A\bx+\bm b$ and $\mathcal Y_i=\phi(\mathcal X_i)=A\mathcal X_i+\bm b$.
Using $\sum_i w_i^m=1$ and the mean-matching property $\sum_i w_i^m\mathcal X_i=\bm m$
(immediate from \eqref{eq:sigma}, since the $\pm$ pairs cancel and $w^m$ is symmetric),
\[
\hat{\bm m}=\sum_i w_i^m(A\mathcal X_i+\bm b)=A\Big(\sum_i w_i^m\mathcal X_i\Big)+\bm b
          =A\bm m+\bm b .
\]
For the covariance, $\mathcal Y_i-\hat{\bm m}=A(\mathcal X_i-\bm m)$, so
\[
\hat{\bm P}=\sum_i w_i^c\,A(\mathcal X_i-\bm m)(\mathcal X_i-\bm m)^\top A^\top
          =A\Big(\sum_i w_i^c(\mathcal X_i-\bm m)(\mathcal X_i-\bm m)^\top\Big)A^\top
          =A\bm P A^\top,
\]
where the last equality is the covariance-matching property of the sigma set: with
$\mathcal X_i-\bm m=\pm[\sqrt{(n+\lambda)\bm P}\,]_j$ for $i=j,\,j+n$ and
$w_j^c=w_{j+n}^c=\tfrac{1}{2(n+\lambda)}$, each pair contributes
$2\cdot\tfrac{1}{2(n+\lambda)}\,(n+\lambda)\,\ell_j\ell_j^\top=\ell_j\ell_j^\top$ with
$\ell_j$ the $j$-th column of $\sqrt{\bm P}$, and $\sum_j\ell_j\ell_j^\top=\bm P$; the
$i=0$ term drops because $\mathcal X_0-\bm m=\bm 0$. Both identities hold for any
$(\alpha,\beta,\kappa)$ giving $n+\lambda\neq0$. \hfill$\qed$

\begin{thebibliography}{9}

\bibitem{julier2004}
S.~J. Julier and J.~K. Uhlmann.
\newblock Unscented filtering and nonlinear estimation.
\newblock \emph{Proceedings of the IEEE}, 92(3):401--422, 2004.
\newblock \href{https://doi.org/10.1109/JPROC.2003.823141}{doi:10.1109/JPROC.2003.823141}.

\bibitem{wan2000}
E.~A. Wan and R.~van der Merwe.
\newblock The unscented Kalman filter for nonlinear estimation.
\newblock In \emph{Proc. IEEE Adaptive Systems for Signal Processing,
Communications, and Control (AS-SPCC)}, pages 153--158, 2000.
\newblock \href{https://doi.org/10.1109/ASSPCC.2000.882463}{doi:10.1109/ASSPCC.2000.882463}.

\bibitem{vandermerwe2004}
R.~van der Merwe.
\newblock \emph{Sigma-Point Kalman Filters for Probabilistic Inference in Dynamic
State-Space Models}.
\newblock PhD thesis, Oregon Health \& Science University, 2004.

\bibitem{sarkka2013}
S.~S\"arkk\"a.
\newblock \emph{Bayesian Filtering and Smoothing}.
\newblock Cambridge University Press, 2013.
\newblock \href{https://doi.org/10.1017/CBO9781139344203}{doi:10.1017/CBO9781139344203}.

\end{thebibliography}

\end{document}
